\newcommand{\appsubsection}[1]{
    \stepcounter{subsection}
    \subsection*{\Alph{section}.\arabic{subsection}\hspace{1em}{#1}}
}

% red, green, and blue comments
\newcommand{\red}[1]{\textcolor{red!80!black}{#1}}
\newcommand{\blue}[1]{\textcolor{blue!80!black}{#1}}
\newcommand{\green}[1]{\textcolor{green!80!black}{#1}}

% highlight in table
\definecolor{sci_light_blue}{HTML}{4869b4}
\definecolor{sci_light_tes}{HTML}{a08aaa}
\definecolor{sci_light_gre}{HTML}{3E5F58}
\definecolor{red_cb}{HTML}{e41a1c}
\definecolor{blue_cb}{HTML}{377eb8}
\definecolor{green_cb}{HTML}{4daf4a}
\definecolor{purple_cb}{HTML}{984ea3}
\definecolor{orange_cb}{HTML}{ff7f00}
\newcommand{\colorr}[1]{\textcolor{red_cb}{#1}}
\newcommand{\colorb}[1]{\textcolor{blue_cb}{#1}}
\newcommand{\colorg}[1]{\textcolor{green_cb}{#1}}
\newcommand{\colorp}[1]{\textcolor{purple_cb}{#1}}
\newcommand{\coloro}[1]{\textcolor{orange_cb}{#1}}

% prettier colors for references
\definecolor{EmeraldGreen}{HTML}{1ea78d}
\definecolor{EnglishRed}{HTML}{b02427}
\hypersetup{colorlinks=true,urlcolor=EmeraldGreen,citecolor=EmeraldGreen,linkcolor=EnglishRed}

% \marginparwidth 4.0cm
% \setlength{\hoffset}{-1.6cm}
% \newcommand{\mpar}[1]{\rule{2pt}{10pt}
%                      {\marginpar{\hbadness10000\flushleft
%                      \sloppy\hfuzz10pt\boldmath\footnotesize#1}}
%                       \typeout{marginpar: #1}\ignorespaces}
% \def\mda{\mpar{\hfil$\downarrow$\hfil}\ignorespaces}
% \def\mua{\mpar{\hfil$\uparrow$\hfil}\ignorespaces}
% \def\mla{\marginpar[\boldmath\hfil$\rightarrow$\hfil]
%                   {\boldmath\hfil$\leftarrow $\hfil}
%                   \typeout{marginpar: $\leftrightarrow$}\ignorespaces}

\newcommand{\eg}{\text{e.g.}\;}
\newcommand{\ie}{\text{i.e.}\;}
\newcommand{\etal}{\textsl{et al}\;}
\newcommand{\eqcomma}{\quad\text{,}} 	% equation comma
\newcommand{\eqperiod}{\quad\text{.}} 	% equation period


% ML-related definitions
\newcommand{\lrgen}{\text{LR}_\text{gen}}
\newcommand{\relu}{\text{ReLU}}
\newcommand{\kl}{D_\text{KL}}
\newcommand{\vpd}{\vec{p}_\text{data}}
\newcommand{\vpmd}{\vec{p}_\text{model}}
\newcommand{\psim}{p_\text{sim}}
\newcommand{\psimr}{p_\text{sim}^\text{r}}
\newcommand{\psimp}{p_\text{sim}^\text{p}}
\newcommand{\pdj}[1]{p_{\text{data},#1}}
\newcommand{\pmdj}[1]{p_{\text{model},#1}}
\newcommand{\pl}{p_\text{latent}}
\newcommand{\pd}{p_\text{data}}
\newcommand{\pcs}{p_{\text{cross section}}}
\newcommand{\pmd}{p_\text{model}}
\newcommand{\dcsi}{d\sigma_{i}}
\newcommand{\dcsixi}{d\sigma_{i\xi}}
\newcommand{\pcsi}{p_{\text{cross section},i}}
\newcommand{\pcsixi}{p_{\text{cross section},i\xi}}
\newcommand{\pmdi}{p_{\text{model}, i}}
\newcommand{\pmdixi}{p_{\text{model}, \xi}}
\newcommand{\punf}{p_\text{unfold}}
\newcommand{\xp}{x_\text{part}}
\newcommand{\xr}{x_\text{reco}}
\newcommand{\pp}{p_\text{sim,p}}
\newcommand{\pr}{p_\text{sim,r}}
\newcommand{\pu}{p_\text{unfold}}
\newcommand{\hpcsi}{\bar{p}^{(i)}_{\text{cross section}}}
\newcommand{\hpmdi}{\bar{p}^{(i)}_\text{model}}
\newcommand{\argmin}{\text{argmin}}
\newcommand{\argmax}{\text{argmax}}
%\newcommand{\softmax}{\text{Softmax}}
%\newcommand{\sigmoid}{\text{Sigmoid}}
\newcommand{\Langle}{\big\langle}
\newcommand{\Rangle}{\big\rangle}
\newcommand{\XLangle}{\Big\langle}
\newcommand{\XRangle}{\Big\rangle}
\newcommand{\XXLangle}{\Bigg\langle}
\newcommand{\XXRangle}{\Bigg\rangle}
\newcommand{\XXXLangle}{\Biggl\langle}
\newcommand{\XXXRangle}{\Biggr\rangle}
%\newcommand{\cbox}[1]{\parbox{4em}{\centering {#1}}}

\newcommand{\ai}{\alpha_i}
\newcommand{\aix}{\alpha_{i \xi}}
%\newcommand{\ajx}{\alpha_{j \xi}}
\newcommand{\dix}{\Delta_{i \xi}}
\newcommand{\Git}{G_{i \theta}}
\newcommand{\git}{g_{i \theta}}
\newcommand{\gbar}{\overline{G}}
\newcommand{\fd}{f_\text{data}}
\newcommand{\mwith}{\text{with}}
\newcommand{\mand}{\text{and}}
\newcommand{\mfor}{\text{for}}

% brackets and parentheses
\newcommand{\Lbrack}{\bigl\lbrack}
\newcommand{\Rbrack}{\bigr\rbrack}
\newcommand{\XLbrack}{\Bigl\lbrack}
\newcommand{\XRbrack}{\Bigr\rbrack}
\newcommand{\XXLbrack}{\biggl\lbrack}
\newcommand{\XXRbrack}{\biggr\rbrack}
\newcommand{\Lbrace}{\bigl\lbrace}
\newcommand{\Rbrace}{\bigr\rbrace}
\newcommand{\XLbrace}{\Bigl\lbrace}
\newcommand{\XRbrace}{\Bigr\rbrace}
\newcommand{\XXLbrace}{\biggl\lbrace}
\newcommand{\XXRbrace}{\biggr\rbrace}
\newcommand{\braces}[1]{\lbrace#1\rbrace}


\newcommand{\qqquad}{\qquad\quad}
\newcommand{\qqqquad}{\qquad\qquad}

% general math definitions
\def\d{\mathrm{d}}
\DeclareMathOperator\erf{erf}
\newcommand\one{\leavevmode\hbox{\small1\normalsize\kern-.33em1}}
\newcommand{\diff}{\mathop{}\!\mathrm{d}} 	% differential
\newcommand{\del}{\partial} 				% partial derivative
\newcommand{\vect}[1]{\bm{#1}} 				% bold vector notation
\newcommand{\abs}[1]{\lvert#1\rvert} 		% absolute value (single vertical lines)
\newcommand{\norm}[1]{\lVert#1\rVert} 		% norm (double vertical lines)
\newcommand{\Z}{\mathbb{Z}} 				% integers
\newcommand{\Q}{\mathbb{Q}} 				% rational numbers
\newcommand{\R}{\mathbb{R}} 				% real numbers
\newcommand{\var}{\operatorname{Var}} 		% variance
\newcommand{\sign}{\operatorname{sign}} 	% sign
\newcommand{\tr}{\operatorname{Tr}}			% trace
\newcommand{\order}{\mathcal{O}} 			% order
\newcommand{\imag}{\mathrm{i}} 				% imaginary unit
\newcommand{\euler}{\mathrm{e}} 			% Euler's number
\newcommand{\iid}{\text{i.i.d.}} 			% independent and identically distributed
\newcommand{\range}[2]{#1,\ldots,#2}
\newcommand{\really}{\stackrel{!}{=}}
\newcommand{\mean}[1]{\left\langle#1\right\rangle}

% particle physics and machine learning definitions
\newcommand{\similarity}{\operatorname{sim}}
\newcommand{\softmax}{\operatorname{Softmax}}
\newcommand{\sigmoid}{\operatorname{Sigmoid}}
\newcommand{\ReLU}{\operatorname{ReLU}}
\newcommand{\LeakyReLU}{\operatorname{LeakyReLU}}
\newcommand{\Attention}{\operatorname{Attention}}
\newcommand{\MultiHead}{\operatorname{MultiHead}}
\newcommand{\Concat}{\operatorname{Concat}}
\newcommand{\head}{\operatorname{head}}

\newcommand{\pT}{p_{\mathrm{T}}} 	% transverse momentum
\newcommand{\pTi}{p_{\mathrm{T},i}} % transverse momentum (with index i)
\newcommand{\pTjet}{p_{\mathrm{T},\,\mathrm{jet}}}
\newcommand{\etajet}{\eta_{\,\mathrm{jet}}}
\newcommand{\yjet}{y_{\,\mathrm{jet}}}
\newcommand{\kT}{k_{\mathrm{T}}}
\newcommand{\mT}{m_{\mathrm{T}}} 	% transverse mass
\newcommand{\sqrts}{\sqrt{s}} 		% center-of-mass energy
\newcommand{\loss}{\mathcal{L}} 	% loss value
\newcommand{\normal}{\mathcal{N}} 	% loss value
\newcommand{\dth}{\Delta_{\mathrm{th/sys}}}
\newcommand{\dst}{\Delta_{\mathrm{stat}}}
\newcommand{\dsy}{\Delta_{\mathrm{syst}}}

\newcommand{\pTa}{p_{\mathrm{T},a}} % transverse momentum (with index a)
\newcommand{\pTb}{p_{\mathrm{T},b}} % transverse momentum (with index b)
\newcommand{\batch}{\mathrm{batch}}
\newcommand{\Nbatch}{N_{\mathrm{batch}}}
\newcommand{\nc}{n_{C}} % number of constituents
\newcommand{\lag}{\mathscr{L}}

\newcommand{\etaphi}{$\eta$--$\phi$\xspace} % eta-phi plane
\newcommand{\AUC}{\mathrm{AUC}} % area under the ROC curve
\newcommand{\es}{\epsilon_{s}} 	% signal efficiency (true-positive rate)
\newcommand{\eb}{\epsilon_{b}} 	% background mistag rate (false-positive rate)
\newcommand{\imtafe}{\epsilon_{b}^{-1}(\epsilon_{s}\!=\!0.5)}



% Software used
\newcommand{\ma}{\textsc{MadAgents}\xspace}
\newcommand{\herwig}{\textsc{Herwig}\xspace}
\newcommand{\vegas}{\textsc{Vegas}\xspace}
\newcommand{\vegaz}{\textsc{VegaZ}\xspace}
\newcommand{\madgraph}{\textsc{MadGraph5\_aMC@NLO}\xspace}
\newcommand{\mg}{\textsc{MadGraph}\xspace}
\newcommand{\pythia}{\textsc{Pythia}\xspace}
\newcommand{\fastjet}{\textsc{FastJet}\xspace}
\newcommand{\tensorflow}{\textsc{TensorFlow}\xspace}
\newcommand{\delphes}{\textsc{Delphes}\xspace}
\newcommand{\pdfset}{\texttt{PDF4LHC15\_nlo\_30}\xspace}
\newcommand{\pytorch}{\textsc{PyTorch}\xspace}
\newcommand{\sherpa}{\textsc{Sherpa}\xspace}
\newcommand{\bonsay}{\textsc{Bonsay}\xspace}
\newcommand{\keras}{\textsc{Keras}\xspace}
\newcommand{\adam}{\textsc{Adam}\xspace}
\newcommand{\madnis}{\textsc{MadNIS}\xspace}
\newcommand{\insane}{\textsc{InSaNe}\xspace}
\newcommand{\lhapdf}{\textsc{Lhapdf~6}\xspace}
\newcommand{\python}{\textsc{Python}\xspace}
\newcommand{\whizard}{\textsc{Whizard}\xspace}
\newcommand{\vamp}{\textsc{Vamp}\xspace}
\renewcommand{\root}{\textsc{Root}\xspace}
%\newcommand{\sklearn}{\texttt{scikit-learn}\xspace}
%\newcommand{\PyTorch}{\texttt{PyTorch}\xspace}
%\newcommand{\Adam}{\texttt{Adam}\xspace}
%\newcommand{\AdamW}{\texttt{AdamW}\xspace}
%\newcommand{\EnergyFlow}{\texttt{EnergyFlow}\xspace}
%\newcommand{\FastJet}{\textsc{FastJet}\xspace}
%\newcommand{\Pythia}{\textsc{Pythia}\xspace}
%\newcommand{\Delphes}{\textsc{Delphes}\xspace}
%\newcommand{\Python}{\texttt{Python}\xspace}
%\newcommand{\NumPy}{\texttt{NumPy}\xspace}
%\newcommand{\Vegas}{\textsc{Vegas}\xspace}
%\newcommand{\Madgraph}{\textsc{Madgraph}\xspace}
%\newcommand{\Sherpa}{\textsc{Sherpa}\xspace}
%\newcommand{\Keras}{\textsc{Keras}\xspace}
%\newcommand{\pysecdec}{py\textsc{SecDec}\xspace}
%\newcommand{\secdec}{\textsc{SecDec}\xspace}

\newcommand{\EFP}{\mathrm{EFP}} 		% energy flow polynomials
\newcommand{\ttbar}{t\bar{t}} 			% top-antitop pair
\newcommand{\electron}{\mathrm{e}^{-}} 	% electron
\newcommand{\positron}{\mathrm{e}^{+}} 	% positron

% hyperlink references
\newcommand{\urlx}[1]{\href{#1}{#1}}
\newcommand{\DOI}[2]{\href{http://dx.doi.org/#2}{#1}}
%\newcommand{\arXiv}[1]{\href{http://arxiv.org/abs/#1}{arXiv:#1}}
\newcommand{\arXiv}[2][]{%
	\ifthenelse{\equal{#1}{}}%
	{\href{http://arxiv.org/abs/#2}{arXiv:#2}}%
	{\href{http://arxiv.org/abs/#2}{arXiv:#2~[#1]}}}

% units of measure
\newcommand{\cmu}{\text{cm}}
\newcommand{\mev}{\text{MeV}}
\newcommand{\gev}{\text{GeV}}
\newcommand{\tev}{\text{TeV}}
\newcommand{\fb}{\text{fb}}
\newcommand{\ab}{\text{ab}}
\newcommand{\pb}{\text{pb}}
\newcommand{\br}{\text{BR}}
\newcommand{\ifb}{\ensuremath{\text{fb}^{-1}} }
\newcommand{\iab}{\ensuremath{\text{ab}^{-1}} }
\newcommand{\ipb}{\text{pb}^{-1}}
\newcommand{\sig}{\text{sig}}
\newcommand{\bkg}{\text{bkg}}

% really great macro by Chris Lester
\def\slashchar#1{\setbox0=\hbox{$#1$}           % set a box for #1
   \dimen0=\wd0                                 % and get its size
   \setbox1=\hbox{/} \dimen1=\wd1               % get size of /
   \ifdim\dimen0>\dimen1                        % #1 is bigger
      \rlap{\hbox to \dimen0{\hfil/\hfil}}      % so center / in box
      #1                                        % and print #1
   \else                                        % / is bigger
      \rlap{\hbox to \dimen1{\hfil$#1$\hfil}}   % so center #1
      /                                         % and print /
   \fi}
\newcommand{\dslash}{\slashchar{\partial}}
\newcommand{\Dslash}{\slashchar{D}}

%\DeclareMathOperator{\tr}{Tr}

% for transformation matrix A
\newcommand{\tikznode}[2]{%
\ifmmode%
\tikz[remember picture,baseline=(#1.base),inner sep=0pt] \node (#1) {$#2$};%
\else
\tikz[remember picture,baseline=(#1.base),inner sep=0pt] \node (#1) {#2};%
\fi}

% Nice looking particle names and masses
\def\mathswitchr#1{\relax\ifmmode{\mathrm{#1}}\else$\mathrm{#1}$\xspace\fi}
\def\mathswitch#1{\relax\ifmmode#1\else$#1$\xspace\fi}
\newcommand{\Pf}{\mathswitch  f}
\newcommand{\Pfbar}{\mathswitch{\bar f}}
\newcommand{\Pq}{\mathswitch q}
\newcommand{\Pqbar}{\mathswitch{\bar q}}
\newcommand{\PW}{\mathswitchr W}
\newcommand{\PWp}{\mathswitchr {W^+}}
\newcommand{\PWm}{\mathswitchr {W^-}}
\newcommand{\PWpm}{\mathswitchr {W^\pm}}
\newcommand{\PZ}{\mathswitchr Z}
\newcommand{\PZp}{\mathswitchr {Z'}}
\newcommand{\Pg}{\mathswitchr g}
\newcommand{\PH}{\mathswitchr H}
\newcommand{\Pe}{\mathswitchr e}
\newcommand{\Pne}{\mathswitch \nu_{\mathrm{e}}}
\newcommand{\Pep}{\mathswitchr {e^+}}
\newcommand{\Pem}{\mathswitchr {e^-}}
\newcommand{\Pmup}{\mu^+}
\newcommand{\Pgngm}{\nu_\mu}
\newcommand{\Pp}{\mathswitchr p}
\newcommand{\Pd}{\mathswitchr d}
\newcommand{\Pu}{\mathswitchr u}
\newcommand{\Ps}{\mathswitchr s}
\newcommand{\Pc}{\mathswitchr c}
\newcommand{\Pb}{\mathswitchr b}
\newcommand{\Pt}{\mathswitchr t}
\newcommand{\Pj}{\mathswitchr j}
\newcommand{\Ptbar}{\mathswitchr{\bar t}}
\newcommand{\Pubar}{\mathswitchr{\bar u}}
\newcommand{\Pdbar}{\mathswitchr{\bar d}}
\newcommand{\jet}{\mathrm{jet}}
\newcommand{\jets}{\mathrm{jets}}
\newcommand{\Pl}{\ell}
\newcommand{\Mf}{\mathswitch {m_\Pf}}
\newcommand{\Ml}{\mathswitch {m_\Pl}}
\newcommand{\Mq}{\mathswitch {m_\Pq}}
\newcommand{\MV}{\mathswitch {M_\PV}}
\newcommand{\MW}{\mathswitch {M_\PW}}
\newcommand{\MA}{\mathswitch {\lambda}}
\newcommand{\MZ}{\mathswitch {M_\PZ}}
\newcommand{\MZp}{\mathswitch {M_\PZp}}
\newcommand{\MH}{\mathswitch {M_\PH}}
\newcommand{\Me}{\mathswitch {m_\Pe}}
\newcommand{\Mmy}{\mathswitch {m_\mu}}
\newcommand{\Mta}{\mathswitch {m_\tau}}
\newcommand{\Md}{\mathswitch {m_\Pd}}
\newcommand{\Mu}{\mathswitch {m_\Pu}}
\newcommand{\Ms}{\mathswitch {m_\Ps}}
\newcommand{\Mc}{\mathswitch {m_\Pc}}
\newcommand{\Mb}{\mathswitch {m_\Pb}}
\newcommand{\Mt}{\mathswitch {m_\Pt}}
\newcommand{\mat}{\mathcal{M}}

\newcommand{\cbox}[1]{\parbox{4em}{\centering {#1}}}

\newcommand{\Ntrain}{\ensuremath{N_\text{train}}\xspace}

\newcommand{\ETmiss}{\ensuremath{E_{T,\text{miss}}}\xspace}

\newcommand{\tilOHD}{\ensuremath{\tilde{\mathcal{O}}_{\Phi D}}\xspace}
\newcommand{\OHD}{\ensuremath{\mathcal{O}_{\Phi D}}\xspace}
\newcommand{\OHbox}{\ensuremath{\mathcal{O}_{\Phi\Box}}\xspace}
\newcommand{\OHW}{\ensuremath{\mathcal{O}_{\Phi W}}\xspace}
\newcommand{\OHB}{\ensuremath{\mathcal{O}_{\Phi B}}\xspace}
\newcommand{\OHWB}{\ensuremath{\mathcal{O}_{\Phi WB}}\xspace}
\newcommand{\OWWW}{\ensuremath{\mathcal{O}_{WWW}}\xspace}
\newcommand{\OHq}{\ensuremath{\mathcal{O}_{\Phi q}^{(3)}}\xspace}
\newcommand{\OHqm}{\ensuremath{\mathcal{O}_{\Phi q}^{(-)}}\xspace}

\newcommand{\cHD}{\ensuremath{c_{\Phi D}}\xspace}
\newcommand{\cHbox}{\ensuremath{c_{\Phi\Box}}\xspace}
\newcommand{\cHW}{\ensuremath{c_{\Phi W}}\xspace}
\newcommand{\cHB}{\ensuremath{c_{\Phi B}}\xspace}
\newcommand{\cHWB}{\ensuremath{c_{\Phi WB}}\xspace}
\newcommand{\cWWW}{\ensuremath{c_{WWW}}\xspace}
\newcommand{\cHq}{\ensuremath{c_{\Phi q}^{(3)}}\xspace}
\newcommand{\cHqm}{\ensuremath{c_{\Phi q}^{(-)}}\xspace}

\newcommand{\HBc}[1]{\textcolor{red}{[hb: #1]}}

% Indent fix for boxes
\makeatletter
\newcommand{\indentafterblock}{%
  \futurelet\pb@next\pb@after
}
\newcommand{\pb@after}{%
  \ifx\pb@next\par
    % blank line next -> keep normal indentation
  \else
    % if we're still in horizontal mode, end the line first
    \ifhmode\par\fi
    % start next paragraph without indentation
    \noindent\ignorespaces
  \fi
}
\makeatother

% Bash, prompt and markdown blocks
\newtcblisting{bashblock}{
  listing engine=minted,
  minted language=bash,
  breakable,
  enhanced,
  colback=black!2,
  colframe=black!20,
  arc=2pt,
  boxrule=0.45pt,
  left=4pt,right=4pt,top=2pt,bottom=2pt,
  boxsep=1pt,
  listing only,
  before skip=4pt, after skip=4pt,
  after=\indentafterblock,
  minted options={
    breaklines,
    breaksymbolleft=,
    fontsize=\footnotesize,
    autogobble,
    tabsize=2
  }
}

\tcbuselibrary{skins,breakable,listings}

\tcbset{
  promptblock/base/.style={
    breakable,
    enhanced,
    arc=2pt,
    boxrule=0.45pt,
    left=4pt,right=12pt,top=4pt,bottom=3pt,
    boxsep=1pt,
    title={#1},
    title filled,
    coltitle=black,
    fontupper=\footnotesize,
    fonttitle=\ttfamily\small\bfseries,
    toptitle=1.5pt,bottomtitle=1.5pt,
    lefttitle=4pt,righttitle=4pt,
    before skip=4pt, after skip=4pt,
    after=\indentafterblock
  }
}

% Bad backup
% \tcbset{
%   promptblock/user/.style={
%     colback=sci_light_blue!5,
%     colframe=sci_light_blue!50,
%     colbacktitle=sci_light_blue!50
%   },
%   promptblock/plan/.style={
%   colback=brown!1.5!yellow!2!white,
%   colframe=brown!6!yellow!8!black,
%   colbacktitle=brown!2!yellow!7!white
% },
%   promptblock/final/.style={
%     colback=green!2!white,
%     colframe=green!20!black,
%     colbacktitle=green!10!gray
%   }
% }

\tcbset{
  promptblock/plan/.style={
    colback=sci_light_tes!5,
    colframe=sci_light_tes!50,
    colbacktitle=sci_light_tes!50
  },
  promptblock/user/.style={
    colback=sci_light_gre!5,
    colframe=sci_light_gre!50,
    colbacktitle=sci_light_gre!50 
},
  promptblock/final/.style={
    colback=sci_light_blue!5,
    colframe=sci_light_blue!50,
    colbacktitle=sci_light_blue!50
  }
}

\newtcolorbox{userblock}[1][]{
  promptblock/base={#1},
  promptblock/user
}

\newtcolorbox{planblock}[1][]{
  promptblock/base={#1},
  promptblock/plan
}

\newtcolorbox{finalblock}[1][]{
  promptblock/base={#1},
  promptblock/final
}

\newtcolorbox{promptblock}[1][]{
  breakable,
  enhanced,
  colback=sci_light_blue!5,
  colframe=sci_light_blue!50,
  arc=2pt,
  boxrule=0.45pt,
  left=4pt,right=12pt,top=4pt,bottom=3pt,
  boxsep=1pt,
  title={#1},
  title filled,
  colbacktitle=sci_light_blue!50,
  coltitle=black,
  fontupper=\footnotesize,
  fonttitle=\ttfamily\small\bfseries,
  toptitle=1.5pt,bottomtitle=1.5pt,
  lefttitle=4pt,righttitle=4pt,
  before skip=4pt, after skip=4pt,
  after=\indentafterblock
}

\newtcblisting{promptblocklit}[1][]{
  breakable,
  enhanced,
  colback=black!3,
  colframe=black!35,
  arc=2pt,
  boxrule=0.45pt,
  left=2pt,right=2pt,top=2pt,bottom=2pt,
  boxsep=0.5pt,
  title={#1},
  title filled,
  colbacktitle=black!20,
  coltitle=black,
  fonttitle=\ttfamily\footnotesize\bfseries,
  toptitle=1pt,bottomtitle=1pt,
  lefttitle=2pt,righttitle=2pt,
  before skip=3pt, after skip=3pt,
  listing only,
  listing engine=listings,
  listing options={
    basicstyle=\ttfamily\tiny\linespread{0.95}\selectfont,
    columns=fullflexible,
    keepspaces=true,
    showspaces=false,
    showstringspaces=false,
    breaklines=true,
    breakatwhitespace=false,
    upquote=true,
    aboveskip=0pt,
    belowskip=0pt,
    literate={—}{{---}}1 {“}{{``}}1 {”}{{''}}1 {→}{{->}}2 {–}{{--}}1 {‘}{{`}}1 {’}{{'}}1 {…}{{...}}3 {√}{{sqrt}}4
  },
  after=\indentafterblock
}


% \usepackage[fencedCode]{markdown} % enables ``` fenced code blocks
% add more options if you want (tables, footnotes, etc.)
\let\circledS\relax
% \usepackage[fencedCode,pipeTables,texMathDollars,tightLists]{markdown}
% \usepackage[fencedCode,pipeTables,texMathDollars,tightLists,finalizeCache,cacheDir=md-cache]{markdown}
\usepackage[fencedCode,pipeTables,texMathDollars,tightLists,frozenCache,cacheDir=md-cache]{markdown}
% \usepackage[fencedCode,pipeTables,texMathDollars,tightLists,finalizeCache,cacheDir=markdown]{markdown}
%\usepackage[fencedCode,pipeTables,texMathDollars,tightLists,frozenCache,cacheDir=markdown]{markdown}

\AddToHook{begindocument/end}{%
  % Force line breaks to behave as spaces
  \def\markdownRendererSoftLineBreak{\unskip\space}%
  \def\markdownRendererHardLineBreak{\unskip\space}%
  % Also force the key-based versions (in case some theme rebinds later)
  \markdownSetup{
    renderers = {
      softLineBreak = {\unskip\space},
      hardLineBreak = {\unskip\space},
    },
  }%
}

\usepackage{fvextra}
\usepackage{xurl}
\usepackage{xstring}

\usetikzlibrary{shadows.blur}

\definecolor{mdAccent}{HTML}{2563EB} % blue-600
\definecolor{mdBg}{HTML}{F8FAFC}     % slate-50
\definecolor{mdFrame}{HTML}{CBD5E1}  % slate-300
\definecolor{mdTitleBg}{HTML}{EFF6FF}% blue-50

\tcbset{
  mdCard/.style={
    enhanced,
    breakable,
    colback=mdBg,
    colframe=mdAccent!35!mdFrame,
    arc=2.2mm,
    boxrule=0.45pt,
    left=1.8mm,right=1.8mm,top=1.3mm,bottom=1.3mm,
    boxsep=0.9mm,
    colbacktitle=mdTitleBg,
    coltitle=black,
    fonttitle=\ttfamily\small\bfseries,
    boxed title style={sharp corners, boxrule=0pt},
    toptitle=0.8mm,bottomtitle=0.8mm,
    before skip=4pt,
    after skip=4pt
  }
}

\newtcolorbox{mdcodeblock}[1][]{%
  enhanced,
  breakable,
  colback=mdBg,
  colframe=mdAccent!30!mdFrame,
  arc=2.0mm,
  boxrule=0.4pt,
  left=4.5pt,right=4.5pt,top=2.2pt,bottom=2.2pt,
  boxsep=1pt,
  #1
}

\newcommand{\mdmaximgwidth}{0.75\linewidth}
\newcommand{\mdmaximgheight}{0.75\textheight}

\makeatletter

\def\markdownRendererHardLineBreakPrototype{\space}
\def\markdownRendererSoftLineBreakPrototype{\space}
\def\markdownRendererSoftLineBreak{\space}
\def\markdownRendererHardLineBreak{\space}


\def\markdownRendererCodeSpan#1{%
  {\ttfamily\footnotesize
    \begingroup
      % Convert markdown escape tokens into literal characters
      \def\markdownRendererUnderscore##1{\string_}%
      \def\markdownRendererBackslash##1{\string\\}%
      \def\markdownRendererLeftBrace##1{\string\{}%
      \def\markdownRendererRightBrace##1{\string\}}%
      \def\markdownRendererDollarSign##1{\string$}%
      \def\markdownRendererTilde##1{\string~}%
      \def\markdownRendererCircumflex##1{\string^}%
      % Expand #1 with the above conversions applied
      \edef\md@codespan{#1}%
      % Use \path so TeX can break long identifiers nicely
      \expandafter\path\expandafter{\md@codespan}%
    \endgroup
  }%
}

% Indented code blocks
\def\markdownRendererInputVerbatim#1{%
  \begin{mdcodeblock}
    \VerbatimInput[
      fontsize=\footnotesize,
      breaklines=true,
      breakanywhere=true
    ]{#1}%
  \end{mdcodeblock}%
}

% Fenced code blocks (```lang)
\def\markdownRendererInputFencedCode#1#2#3{%
  \begin{mdcodeblock}
    \VerbatimInput[%
fontsize=\footnotesize,%
breaklines=true,%
breakanywhere=true%
    ]{#1}%
  \end{mdcodeblock}%
}

\def\markdownRendererImage#1#2#3{%
  \begin{center}
    \includegraphics[width=\mdmaximgwidth,height=\mdmaximgheight,keepaspectratio]{#3}%
  \end{center}%
}

\makeatother

\newcommand{\markdownfile}[2][Markdown File]{%
  \begin{tcolorbox}[
    mdCard,
    title={#1},
    fontupper=\footnotesize,
    after=\indentafterblock
  ]
    \begingroup
      % No ToC entries from inside the Markdown
      \renewcommand\addcontentsline[3]{}%

      % Avoid extra paragraph glue inside the box
      \setlength{\parskip}{0pt}%
      \setlength{\parindent}{0pt}%

      % Make Markdown headings small and NOT affect counters
      \renewcommand{\section}[1]{\par\smallskip\noindent\textbf{\footnotesize ##1}\par\smallskip}
      \renewcommand{\subsection}[1]{\par\smallskip\noindent\textbf{\footnotesize ##1}\par\smallskip}
      \renewcommand{\subsubsection}[1]{\par\smallskip\noindent\textbf{\footnotesize ##1}\par\smallskip}
      \renewcommand{\paragraph}[1]{\par\smallskip\noindent\textbf{\footnotesize ##1}\par\smallskip}
      \renewcommand{\subparagraph}[1]{\par\smallskip\noindent\textbf{\footnotesize ##1}\par\smallskip}

      % Tighten list spacing
      \setlist[itemize,1]{topsep=0pt,itemsep=0pt,parsep=0pt,partopsep=0pt}
      \setlist[itemize,2]{topsep=0pt,itemsep=0pt,parsep=0pt,partopsep=0pt}
      \setlist[enumerate,1]{topsep=0pt,itemsep=0pt,parsep=0pt,partopsep=0pt}
      \setlist[enumerate,2]{topsep=0pt,itemsep=0pt,parsep=0pt,partopsep=0pt}

      \markdownInput{#2}%
      
      \endgroup
  \end{tcolorbox}%
}
